\documentclass[12pt, letterpaper]{article}

\usepackage[utf8]{inputenc}
\usepackage[T1]{fontenc}
\usepackage[spanish]{babel}
\usepackage{geometry}
\usepackage{amsmath, amssymb, amsthm}
\usepackage{booktabs}
\usepackage{array}
\usepackage{microtype}
\usepackage{setspace}
\usepackage{parskip}
\usepackage{titlesec}
\usepackage{fancyhdr}
\usepackage{enumitem}
\usepackage{bm}
\usepackage{mathtools}
\usepackage{tcolorbox}
\usepackage{hyperref}

\tcbuselibrary{skins, breakable}

\geometry{top=2.8cm, bottom=3.0cm, left=3.2cm, right=3.2cm}
\setstretch{1.20}

\pagestyle{fancy}
\fancyhf{}
\fancyhead[C]{\small\textsc{Econometria · Gujarati --- Ejercicio 10.9}}
\fancyfoot[C]{\small\thepage}
\renewcommand{\headrulewidth}{0.4pt}
\renewcommand{\footrulewidth}{0pt}

\titleformat{\section}
  {\large\bfseries\scshape}{\thesection.}{0.6em}{}
  [\vspace{-4pt}\rule{\linewidth}{0.4pt}]
\titleformat{\subsection}{\normalsize\bfseries}{\thesubsection.}{0.5em}{}
\titleformat{\subsubsection}{\normalsize\itshape}{}{0em}{}

\newtheorem{prop}{Proposicion}
\newtheorem{lema}{Lema}
\theoremstyle{remark}
\newtheorem{obs}{Observacion}

\newtcolorbox{clave}{
  colback=white, colframe=black,
  boxrule=0.5pt, arc=3pt,
  left=8pt, right=8pt, top=6pt, bottom=6pt,
  breakable
}

\newcommand{\E}{\mathbb{E}}
\newcommand{\Var}{\operatorname{Var}}
\newcommand{\Cov}{\operatorname{Cov}}
\newcommand{\plim}{\operatorname{plim}}

\begin{document}

\begin{center}
  {\LARGE\bfseries Funcion de Produccion Cobb-Douglas,}\\[6pt]
  {\LARGE\bfseries Multicolinealidad y Sesgo de Especificacion}\\[12pt]
  {\large\itshape Sector manufacturero de los 50 estados de EE.UU., 2005}\\[14pt]
  \rule{0.55\linewidth}{0.6pt}\\[8pt]
  {\normalsize Ejercicio~10.9 --- \textit{Econometria}, Gujarati \& Porter, 5.a ed.}\\[4pt]
  {\small\textsc{Correlacion entre insumos, decision de omision y naturaleza del sesgo resultante}}
\end{center}

\vspace{10pt}

\noindent\textbf{Contexto.}\quad
El ejemplo del capitulo~7 ajusta una funcion de produccion Cobb-Douglas para
el sector manufacturero de los 50 estados y el Distrito de Columbia de
Estados Unidos en 2005. En su forma logaritmica lineal, el modelo es:

\begin{equation}
  \ln Q_i = \beta_1 + \beta_2 \ln L_i + \beta_3 \ln K_i + u_i
  \label{eq:cobb_douglas}
\end{equation}

\noindent donde $Q_i$ es el valor del producto manufacturero del estado $i$,
$L_i$ es el numero de empleados (insumo trabajo) y $K_i$ es el stock de
capital utilizado. Los resultados reportados en~(7.9.4) muestran que tanto
$\hat\beta_2$ como $\hat\beta_3$ son estadisticamente significativos de forma
individual.

\section{El modelo Cobb-Douglas y sus supuestos}

\subsection{Forma funcional y propiedades}

La funcion Cobb-Douglas en su forma general es:
\begin{equation}
  Q_i = A \cdot L_i^{\beta_2} \cdot K_i^{\beta_3} \cdot e^{u_i}
  \label{eq:cd_nivel}
\end{equation}
Tomando logaritmos naturales se obtiene el modelo lineal~(\ref{eq:cobb_douglas}).
Los parametros tienen interpretacion directa:

\begin{itemize}[leftmargin=1.6em, itemsep=4pt]
  \item $\beta_2 = \partial \ln Q / \partial \ln L$: elasticidad de la produccion
    respecto al trabajo (incremento porcentual en $Q$ ante un incremento del
    $1\%$ en $L$, manteniendo $K$ constante).
  \item $\beta_3 = \partial \ln Q / \partial \ln K$: elasticidad de la produccion
    respecto al capital.
  \item $\beta_2 + \beta_3$: rendimientos a escala. Si $> 1$, rendimientos
    crecientes; $= 1$, constantes; $< 1$, decrecientes.
\end{itemize}

\subsection{Estructura de los datos de corte transversal estatal}

Los datos son de corte transversal ($n = 51$): una observacion por estado.
A diferencia de las series de tiempo, aqui la variacion proviene de diferencias
\emph{estructurales} entre estados: algunos tienen sectores manufactureros
grandes (Ohio, Michigan, Texas) y otros pequenos. Esta variacion es la que
permite identificar los parametros, pero tambien genera la correlacion entre
$\ln L$ y $\ln K$ que se analiza a continuacion.

\section{Inciso (a): correlacion entre trabajo y capital}

\subsection{Por que se espera alta correlacion}

En datos de produccion por unidad geografica (estados, paises, industrias),
el trabajo y el capital tienden a moverse juntos por tres razones:

\begin{enumerate}[leftmargin=1.8em, itemsep=6pt]
  \item \textbf{Efecto escala}: los estados con mayor actividad manufacturera
    emplean mas trabajadores \emph{y} mas capital simultaneamente. Un estado
    con un sector manufacturero diez veces mayor que otro tendera a tener
    mas de ambos insumos. La variacion en el tamano del sector domina la
    variacion en las intensidades relativas de uso de los factores.

  \item \textbf{Complementariedad tecnologica}: en muchos procesos industriales
    el capital (maquinaria, equipos) requiere trabajadores para operarlo.
    La relacion capital-trabajo $K/L$ es relativamente estable dentro de
    una misma industria, lo que implica que $\ln K$ y $\ln L$ son
    aproximadamente proporcionales:
    \begin{equation}
      \ln K_i \approx c + \ln L_i + \varepsilon_i
      \label{eq:proporc}
    \end{equation}
    con $\varepsilon_i$ de varianza pequena relativa a la varianza de
    $\ln L_i$.

  \item \textbf{Homogeneidad industrial}: los 51 estados pertenecen al mismo
    pais con acceso a tecnologias similares. La funcion de produccion es
    aproximadamente la misma en todos. Esto implica relaciones
    capital-trabajo parecidas entre estados, lo que refuerza la correlacion.
\end{enumerate}

\subsection{Medicion de la correlacion}

El coeficiente de correlacion entre los logaritmos de los insumos es:
\begin{equation}
  r_{LK} = \frac{\sum (\ln L_i - \overline{\ln L})(\ln K_i - \overline{\ln K})}
                {\sqrt{\sum (\ln L_i - \overline{\ln L})^2
                       \sum (\ln K_i - \overline{\ln K})^2}}
  \label{eq:corr_LK}
\end{equation}

En aplicaciones tipicas de funciones de produccion con datos de corte transversal,
$r_{LK}$ suele estar en el rango $0.85$--$0.98$. El VIF correspondiente es:
\begin{equation}
  \text{VIF}(\ln L) = \text{VIF}(\ln K) = \frac{1}{1 - r_{LK}^2}
  \label{eq:VIF_LK}
\end{equation}

Para $r_{LK} = 0.95$: $\text{VIF} = 1/(1 - 0.9025) = 10.3$, que ya supera
el umbral convencional de colinealidad severa. Para $r_{LK} = 0.99$:
$\text{VIF} = 50$.

\subsection{La paradoja de la significancia individual}

A pesar de la alta correlacion esperada, el enunciado indica que ambos
coeficientes son estadisticamente significativos. Esto es posible
—aunque infrecuente bajo colinealidad severa— cuando:

\begin{enumerate}[leftmargin=1.8em, itemsep=4pt]
  \item El tamano muestral ($n = 51$) es suficientemente grande para que la
    variacion muestral de $\ln L$ con $\ln K$ aproximadamente constante
    sea informativa.
  \item Los errores estandar, aunque inflados por el VIF, no son tan grandes
    como para borrar la significancia de coeficientes que son economicamente
    grandes (elasticidades del orden de $0.5$--$0.7$).
\end{enumerate}

La relacion entre el estadistico $t$ y el VIF es:
\begin{equation}
  t_j = \frac{\hat\beta_j}{\hat\sigma / \sqrt{S_{jj}}} \cdot \frac{1}{\sqrt{\text{VIF}_j}}
  \label{eq:t_vif}
\end{equation}
Con elasticidades grandes y varianza residual razonable, el $t$ puede ser
significativo incluso con VIF$> 10$.

\begin{clave}
  La respuesta al inciso~(a) es: \textbf{si, se espera alta correlacion} entre
  $\ln L$ y $\ln K$ en datos de produccion estatal, por efecto escala y
  complementariedad tecnologica. Sin embargo, la significancia individual de
  ambos coeficientes indica que los datos tienen suficiente variacion
  independiente para identificar los efectos separados, aunque con precision
  reducida respecto al caso ideal de ortogonalidad.
\end{clave}

\section{Inciso (b): ¿eliminar la variable trabajo?}

\subsection{Respuesta}

\textbf{No}. Eliminar $\ln L$ del modelo es incorrecto por razones teoricas
y econometricas que se detallan a continuacion.

\subsection{Razon 1: fundamento teorico inamovible}

La funcion de produccion Cobb-Douglas postula que la produccion depende
\emph{conjuntamente} del trabajo y el capital:
\begin{equation}
  Q = f(L, K)
  \label{eq:produccion}
\end{equation}
Esta especificacion no es una convencion estadistica: refleja la realidad
tecnologica de que el producto surge de la combinacion de ambos factores.
Eliminar el trabajo equivale a asumir $\beta_2 = 0$, es decir, que los
trabajadores no contribuyen a la produccion. Esto contradice la evidencia
empirica universal.

\subsection{Razon 2: los coeficientes son individualmente significativos}

El criterio correcto para eliminar una variable de un modelo es la
\emph{irrelevancia estadistica y teorica}, no la colinealidad. La presencia
de multicolinealidad no es una razon valida para omitir variables relevantes.
Gujarati (2010) es explicito en este punto: la multicolinealidad es un
``problema de los datos'', no del modelo; el remedio no es alterar la
especificacion teorica.

Si $\hat\beta_2$ es significativo, el test $t$ rechaza $H_0: \beta_2 = 0$,
lo que significa que los datos contienen evidencia de que el trabajo tiene
un efecto sobre la produccion separable del capital. Ignorar esta evidencia
es una decision metodologica incorrecta.

\subsection{Razon 3: costo de la omision (sesgo)}

Omitir $\ln L$ introduce un sesgo sistematico en $\hat\beta_3$. Este sesgo
se deriva en detalle en el inciso~(c).

\begin{clave}
  La respuesta al inciso~(b) es: \textbf{no se debe eliminar el trabajo}.
  La colinealidad entre regresores relevantes no justifica la omision de
  ninguno de ellos. El modelo debe mantenerse completo porque ambos insumos
  son teoricamente esenciales y empiricamente significativos.
\end{clave}

\section{Inciso (c): naturaleza del sesgo de especificacion}

\subsection{Tipo de sesgo}

Al estimar el modelo reducido:
\begin{equation}
  \ln Q_i = \alpha_1 + \alpha_3 \ln K_i + v_i
  \label{eq:reducido}
\end{equation}
omitiendo $\ln L_i$, se incurre en \textbf{sesgo por variable omitida}
(specification bias u omitted variable bias). Este es uno de los sesgos
de especificacion catalogados por Ramsey (1969) y discutidos en el
capitulo~13 de Gujarati.

\subsection{Derivacion formal del sesgo}

El verdadero modelo es~(\ref{eq:cobb_douglas}). Sea $b_{LK}$ el coeficiente
de la regresion auxiliar de $\ln L$ sobre $\ln K$:
\begin{equation}
  b_{LK} = \frac{\sum \ell_i k_i}{\sum k_i^2}
  \label{eq:b_LK}
\end{equation}
donde $\ell_i = \ln L_i - \overline{\ln L}$ y $k_i = \ln K_i - \overline{\ln K}$.

El estimador MCO $\hat\alpha_3$ del modelo reducido~(\ref{eq:reducido}) es:
\begin{equation}
  \hat\alpha_3 = \frac{\sum k_i \ln Q_i}{\sum k_i^2}
  \label{eq:alpha3}
\end{equation}

Sustituyendo el verdadero modelo en~(\ref{eq:alpha3}):
\begin{align}
  \hat\alpha_3
  &= \frac{\sum k_i (\beta_1 + \beta_2 \ln L_i + \beta_3 \ln K_i + u_i)}
          {\sum k_i^2}
  \notag\\[4pt]
  &= \beta_2 \cdot \frac{\sum k_i \ell_i}{\sum k_i^2}
   + \beta_3 \cdot \frac{\sum k_i^2}{\sum k_i^2}
   + \frac{\sum k_i u_i}{\sum k_i^2}
  \notag\\[4pt]
  &= \beta_2 b_{LK} + \beta_3 + \frac{\sum k_i u_i}{\sum k_i^2}
  \label{eq:alpha3_descomp}
\end{align}

Tomando esperanza (bajo exogeneidad de $\ln K$):
\begin{equation}
  \boxed{
    \E[\hat\alpha_3] = \beta_3 + \underbrace{\beta_2 \cdot b_{LK}}_{\text{sesgo}}
  }
  \label{eq:sesgo}
\end{equation}

\subsection{Signo y magnitud del sesgo}

Para determinar el signo del sesgo $\beta_2 \cdot b_{LK}$:

\begin{enumerate}[leftmargin=1.8em, itemsep=6pt]
  \item \textbf{Signo de $\beta_2$}: la elasticidad del producto respecto
    al trabajo es positiva ($\beta_2 > 0$): mas trabajadores implica mas
    produccion, manteniendo el capital constante. En aplicaciones empiricas
    tipicas, $\hat\beta_2 \approx 0.5$--$0.7$.

  \item \textbf{Signo de $b_{LK}$}: el coeficiente de la regresion de
    $\ln L$ sobre $\ln K$ es positivo ($b_{LK} > 0$) porque los estados
    con mas capital tambien tienen mas trabajadores (efecto escala). Dado
    que la correlacion $r_{LK} > 0$, se tiene:
    \begin{equation}
      b_{LK} = r_{LK} \cdot \frac{s_{\ln L}}{s_{\ln K}} > 0
      \label{eq:b_LK_corr}
    \end{equation}
    donde $s_{\ln L}$ y $s_{\ln K}$ son las desviaciones estandar muestrales.
\end{enumerate}

Por tanto, el sesgo es:
\begin{equation}
  \text{Sesgo} = \beta_2 \cdot b_{LK} > 0
  \label{eq:sesgo_positivo}
\end{equation}

El estimador del coeficiente del capital en el modelo reducido \textbf{sobreestima}
la verdadera elasticidad del capital:
\begin{equation}
  \E[\hat\alpha_3] = \beta_3 + \beta_2 b_{LK} > \beta_3
  \label{eq:sobreestimacion}
\end{equation}

\subsection{Interpretacion economica del sesgo}

Cuando omitimos $\ln L$, el coeficiente del capital $\hat\alpha_3$
``absorbe'' parte del efecto del trabajo. Al aumentar $\ln K$, los
datos muestran que $\ln L$ tambien aumenta (porque $b_{LK} > 0$);
ese incremento en $\ln L$ eleva la produccion con elasticidad $\beta_2$.
Pero el modelo reducido no puede separar estas dos contribuciones y
se las atribuye enteramente al capital. Por eso $\hat\alpha_3 > \beta_3$.

\subsection{Ilustracion numerica}

Suponga que los verdaderos parametros y la regresion auxiliar son:
\begin{align}
  \hat\beta_2 &= 0.60 \quad \text{(elasticidad del trabajo)}\\
  \hat\beta_3 &= 0.40 \quad \text{(elasticidad del capital)}\\
  b_{LK}      &= 0.80 \quad \text{(pendiente de } \ln L \text{ sobre } \ln K\text{)}
\end{align}

El sesgo es:
\begin{equation}
  \text{Sesgo} = 0.60 \times 0.80 = 0.48
  \label{eq:sesgo_num}
\end{equation}

El modelo reducido estimaria:
\begin{equation}
  \E[\hat\alpha_3] = 0.40 + 0.48 = 0.88
  \label{eq:alpha3_num}
\end{equation}

Es decir, el coeficiente del capital casi se duplica por el sesgo: en vez
de $0.40$, se obtendria $0.88$. Esta distorsion conduciria a conclusiones
erroneas sobre la productividad marginal del capital y los rendimientos
a escala del sector manufacturero.

\subsection{Consecuencias adicionales sobre la inferencia}

Ademas del sesgo en el coeficiente, omitir $\ln L$ tiene consecuencias
para la varianza estimada del error:

\begin{equation}
  \sigma_v^2 = \Var(u_i + \beta_2 \ln L_i) = \sigma_u^2 + \beta_2^2 \sigma_{\ln L}^2
  \label{eq:var_error_reducido}
\end{equation}

La varianza del error del modelo reducido es mayor que la del completo
($\sigma_v^2 > \sigma_u^2$), lo que infla los errores estandar de $\hat\alpha_3$
y deteriora la precision de la inferencia. Las pruebas $t$ y $F$ del modelo
reducido no son validas para el verdadero parametro $\beta_3$.

\begin{clave}
  \textbf{Sesgo de especificacion — resumen:}
  \begin{enumerate}[leftmargin=1.4em, itemsep=4pt]
    \item \textbf{Tipo}: sesgo por variable omitida (subajuste del modelo).
    \item \textbf{Formula}: $\E[\hat\alpha_3] = \beta_3 + \beta_2 b_{LK}$.
    \item \textbf{Signo}: positivo, porque $\beta_2 > 0$ y $b_{LK} > 0$.
      El coeficiente del capital se sobreestima.
    \item \textbf{Magnitud}: proporcional al efecto real del trabajo omitido
      ($\beta_2$) y a la correlacion entre capital y trabajo ($b_{LK}$).
    \item \textbf{Consecuencia adicional}: la varianza del error del modelo
      reducido es mayor, deteriorando la precision de toda la inferencia.
  \end{enumerate}
\end{clave}

\section{Diagnostico y conclusion practica}

El ejercicio ilustra la tension entre dos problemas:

\begin{itemize}[leftmargin=1.6em, itemsep=5pt]
  \item \textbf{Mantener el modelo completo}: estimadores insesgados pero
    con errores estandar inflados por la colinealidad entre $\ln L$ y $\ln K$.
    Sin embargo, como ambos coeficientes son significativos, la colinealidad
    no impide la inferencia en este caso.

  \item \textbf{Reducir el modelo}: errores estandar menores para el
    coeficiente del capital, pero el estimador es sesgado. El sesgo
    $\beta_2 b_{LK} \approx 0.48$ es mucho mayor que cualquier reduccion
    razonable del error estandar. La precision ganada no compensa el sesgo
    introducido.
\end{itemize}

La decision correcta es mantener el modelo completo~(\ref{eq:cobb_douglas}).
Si la multicolinealidad genera errores estandar inaceptablemente grandes,
los remedios adecuados son agregar datos, imponer restricciones teoricas
(como $\beta_2 + \beta_3 = 1$ para rendimientos constantes a escala) o
usar estimadores con sesgo controlado como la regresion Ridge, no omitir
variables teoricamente relevantes.

\appendix

\section{Condicion de Ramsey para sesgo de especificacion}

Ramsey (1969) mostro que el sesgo de especificacion por omision de variables
relevantes puede detectarse con el test RESET (\textit{Regression Equation
Specification Error Test}). El procedimiento es:

\begin{enumerate}[leftmargin=1.8em, itemsep=4pt]
  \item Estimar el modelo reducido y obtener los valores ajustados $\hat Q_i$.
  \item Regresar los residuos $\hat v_i$ sobre potencias de $\hat Q_i$:
    $\hat v_i = \gamma_1 \hat Q_i^2 + \gamma_2 \hat Q_i^3 + \epsilon_i$.
  \item Contrastar $H_0: \gamma_1 = \gamma_2 = 0$ con un test $F$.
    Un rechazo indica mala especificacion del modelo.
\end{enumerate}

\section{Restriccion de rendimientos constantes a escala}

Si se impone la restriccion teorica $\beta_2 + \beta_3 = 1$
(rendimientos constantes a escala), el modelo~(\ref{eq:cobb_douglas}) puede
reescribirse como:
\begin{equation}
  \ln Q_i - \ln K_i = \beta_1 + \beta_2(\ln L_i - \ln K_i) + u_i
  \label{eq:rcs}
\end{equation}

En~(\ref{eq:rcs}) el unico regresor es $\ln(L_i/K_i)$, la
intensidad de trabajo. Este regresor no sufre la colinealidad del par
$(\ln L, \ln K)$ porque es la \emph{razon} entre ellos, y su variacion
no esta dominada por el efecto escala. La restriccion de rendimientos
constantes elimina simultaneamente la colinealidad y el sesgo de omision,
al costo de imponer una restriccion que debe ser contrastada empiricamente.

\vspace{14pt}
\noindent\rule{\linewidth}{0.4pt}

\noindent{\small\textit{Nota.} Los valores numericos de la ilustracion
son hipoteticos y sirven solo para cuantificar la magnitud tipica del sesgo.
Los valores exactos dependen de los datos del capitulo~7 de Gujarati~(2010).
La formula del sesgo $\E[\hat\alpha_3] = \beta_3 + \beta_2 b_{LK}$ es exacta
en muestras finitas bajo los supuestos clasicos de MCO.}

\end{document}
